\section{Limitations and Possible Improvements}
\subsection{Data Collection}
If the data for each class were obtained from a set of real buildings in the region instead of just one, it would yield better results. Ideally, a statistically significant set of buildings and variable utilization profile could be used. Moreover, it would be beneficial to weigh the simulation data on the real energy requirements of those buildings. This is a much more complex task since it would require extensive data acquisition throughout the year.

From the simulation, humidification and dehumidification electricity needs are neglected since these plants are rarely installed today. A better evaluation would consider the possible future share of air treatment plants in the residential sector. For the purpose of this study, the share of installed air conditioning systems for cooling was assumed constant throughout the years, but it might increase in the future.

Regarding the renewable energy sources, solar-thermal systems were neglected since they are considered not very impactful and complex to evaluate.

Time-dependent efficiency for heat pumps should be accounted for, taking into account the ambient temperature could improve the model.

Assessing appliances electricity demand is the most complex contribution to evaluate, especially considering that it is not easy to forecast how electricity consumption will evolve in 60 years. Collecting a statistically significant number of data from the region regarding the type of appliances installed and their consumption would improve the simulation.
For the purpose of this study the appliances demand was assumed constant throughout the years.

\subsection{Plant Modelling and Simulation}
The Modelica model could be utilized to analyze the dynamic response of the plant to transient conditions, thereby evaluating the system's adequacy and frequency control capabilities. This analysis could focus on critical periods, such as specific days of the year characterized by significant demand oscillations.

Additionally, buffers with limited capacity could be modeled to further understand their impact on system dynamics. The hydrogen plant could be simulated to partially shut down some HTSE units, allowing for better adaptation to demand variations.

Incorporating chemical kinetics into the methanation model could enhance the simulation of the plant's dynamic response to transient conditions. Furthermore, the optimization of hydrogen tapped for market consumption could be based on a detailed analysis of the transportation sector.

Lastly, the sizing of the entire plant should be optimized based on a comprehensive economic assessment.

\subsection{Economic Assessment}


\subsection{Environmental Impact Assessment}
The impact on water resources — primarily needed for reactor operations, Balance of Plant (BoP) systems, and HTSE units — is considered negligible. 
This is largely due to the plant's proximity to Lake Maggiore, Italy's second-largest lake, which provides a reliable and abundant water supply without significantly affecting the its
hydrological balance or ecosystem.

However, a quantitative assessment of the water balance is recommended to determine precise water consumption and discharge rates. This evaluation should also consider potential 
thermal impacts on the lake caused by plant operations. Given the lake's considerable size, such impacts are expected to be minimal, but detailed modeling is necessary for confirmation. 
The results should be compared with applicable regulatory standards for industrial water use.

For modeling purposes, the water required by the electrolysers is assumed to be unlimited, based on the lake's ample supply. Consequently, water costs are excluded from the HTSEs' 
operational expenditure (OPEX), influencing the overall economic analysis.
